\subsection{The two-prime class quotient persists after every finite enlargement}

\begin{statement}{A.153}{Proved}
{Persistence with the actual cyclotomic residue-degree maps}
Keep the fields \(k,F,L_n\) and the integer \(n_0\) of A.151.
Let \(T\) be any fixed finite set of rational places containing
\(2,3,11,59,\infty\), and let \(T_n\) be its inverse image in \(L_n\).
Put
\[
 n_T=\max\left(\{n_0\}\cup
       \{1+v_3(p^2-1):p\in T\setminus\{3,\infty\}\}\right).
\tag{A.153.1}
\]
For \(n\geq n_T\), the same field \(FL_n/L_n\) is a
nontrivial unramified cyclic cubic extension split at every place of
\(T_n\).  Consequently its actual Artin characters give surjections
\[
 \begin{gathered}
 \operatorname{Cl}(\mathcal O_{L_n,T_n})\{3\}
                       \longrightarrow\operatorname{Gal}(F/k)=\Z/3,\\
 \varprojlim_{n\geq n_T}
       \operatorname{Cl}(\mathcal O_{L_n,T_n})\{3\}
                       \longrightarrow\Z/3,
 \end{gathered}
\tag{A.153.2}
\]
where every transition in the inverse limit is ideal norm.
Thus no fixed finite enlargement of the allowed support makes the
three-primary class groups of these Galois layers vanish eventually.
This does not give a nonzero value of the target norm obstruction.
\end{statement}
\begin{proof}
A.151 already proves the assertion at \(2,3,11,59,\infty\),
as well as \(F\cap L_n=k\).
Consider a new rational prime \(p\in T\), and a place \(v\) of
\(k\) above it.  The extension \(F/k\) is unramified at \(v\);
its local degree is either one or three.
For \(p\ne3\), the cyclotomic extension
\(\Q_p(\mu_{3^n})/\Q_p\) is unramified, of degree
\(\operatorname{ord}_{3^n}(p)\).
If \(p\equiv1\pmod3\), this order has three-part
\(3^{\max(0,n-v_3(p-1))}\).
If \(p\equiv-1\pmod3\), its order is
\(2\cdot3^{\max(0,n-v_3(p^2-1))}\).
These identities follow by applying the binomial expansion successively
to \(p\), respectively \(p^2\), in \(1+3\Z_3\).
Since \([k_v:\Q_p]\) is respectively one or two, (A.153.1) makes
the residue degree of
\(k_v(\mu_{3^n})/k_v\) divisible by three.

Inside a fixed local algebraic closure, the unramified extension of
each finite degree is unique.  Hence the unramified cubic local
extension supplied by \(F\), when nontrivial, lies in
\(k_v(\mu_{3^n})\), and therefore in every completion of \(L_n\)
above \(v\).  This is an actual local-field inclusion; it proves that
\(FL_n/L_n\) splits at every new place of \(T_n\).

The Artin maps therefore kill both principal ideals and all primes
inverted in \(\mathcal O_{L_n,T_n}\).
Their restriction to the literal common field \(F\) is unchanged.
The Frobenius identity (A.151.9), or the norm--Artin diagram of
\cite[V, Proposition 3.3]{MilneCFT} with the abelian compositum
specified there, proves their norm compatibility.
For each fixed element of \(\Z/3\), the corresponding finite class
fibres are nonempty.  A point at a largest finite height supplies all
lower points by norm, so compactness of the product of those finite
fibres proves the second surjection in (A.153.2).
\end{proof}

The exact residue-order checks are retained in
\path{code/two_prime_finite_support_persistence_certificate.py}
and its output.  They certify the finite arithmetic used in
(A.153.1); the proof supplies the class-field and inverse-limit maps.
